%#!/usr/bin/pdflatex \catcode35=14 \input
%\catcode`\#=6

\documentclass{beamer}
%\usetheme{Default}
\usepackage{verbatim,subcaption,listings}
\usepackage[format=plain]{caption}

\iffalse 
Sections
  Chapter 1
    Who we are, what is vicharak
    Moores law + dennard scaling + von neumann bottleneck
    plea (new, creative, from-first-principles architectures, languages and 
      compute infra) should be created
    overview of existing compute
    4 problems in one slide
    Key ideas (reconfigurability + Heterogeneity (complement not replace))
    Two problems (problem 1 and 2)
  Chapter 2
    Problem 1 - Lack of programming model
    how to reconfigure
    flow based computation
    exemplary dsl + accompanying image
  Chapter 3
    Problem 2 - verilog is hard 
      verilog is the reference line (problems above/below)
      problems above involve simulation/hdl abstractions/hls
      problems below involve synthesis, placement, routing (faster, better)
    we want to solve the problems below verilog 
      cad flow (i.e. operation below verilog) 
    optimization oppurtunities
  Chapter 4
    realizing this goal
    gati
    periplex
    demo
  Conclusion
\fi

\iffalse

Abstract

As Moore's Law's imminent slowdown takes effect, there is a need for innovative
architectures designed from the ground up with both hardware and software in
mind. We describe an architecture along with a family of software programs
(compilers and runtimes) that form the basis of this hardware-software system.
Interoperation with existing compute infrastructure is proposed. The
architecture executes expected functionality not with a fixed instruction set
but by reconfiguring flexible reprogrammable hardware dynamically at computation
time. This reconfiguration, driven by a software compiler, specializes the
hardware for every function but generalizes through reconfiguration. In contrast
to traditional CPU-memory back-and-forth dataflow, this architecture exhibits a
unilateral flow of data with zero instruction overhead, alleviating the von
Neumann bottleneck. Opportunities to improve existing compilers (EDA Software) for
reconfigurable chips (FPGAs) are also discussed.
\fi

\lstset{language=C++, basicstyle=\ttfamily,
                keywordstyle=\color{blue}\ttfamily,
                stringstyle=\color{red}\ttfamily,
                commentstyle=\color{green}\ttfamily,
                tabsize=1,breaklines=true,numbers=left,
                morecomment=[l][\color{magenta}]{\#}}


\title{No-ISA is the Best ISA}
\subtitle{}
\author{Shreeyash Pandey, Rishik Ram Jallarapu}
\institute{Vicharak, India @ vicharak.in}
\date{28th September, 2024}

%\usebackgroundtemplate%
%{%
%    \includegraphics[width=\paperwidth,height=\paperheight]{logo.eps}%
%}

\usepackage[backend=bibtex,style=numeric]{biblatex}
\addbibresource{slides.bib}

\titlegraphic{
  \includegraphics[width=2cm]{images/logo.png}
}


\begin{document}

\begin{frame}
\titlepage
\end{frame}

{
\logo{\includegraphics[width=.07\textwidth]{images/logo.png}}

\begin{frame}[fragile]
\frametitle{About us}

  Goals:
  \begin{itemize}
    \item Introduce software-level reconfigurability to hardwares
    \item Develop consumer and industrial computing hardware for a new ecosystem: desktop, mobile, and cloud-compatible
  \end{itemize}

Vicharak's team of ~50 engineers and thinkers is committed to realizing this vision across all verticals.

  - Akshar Vastarpara (Founder and CEO, Vicharak)

\framesubtitle{}
\end{frame}

\begin{frame}[fragile]
\frametitle{Contents}

  \begin{itemize}
    \item Chapter 1 - Motivations for our Work
    \item Chapter 2 - Introduction To Reconfigurable And Heterogeneous Computing
    \item Chapter 3 - Need for modern EDA Compilers
    \item Chapter 4 - Work Done Towards Implementation  
  \end{itemize}
\end{frame}

\begin{frame}[c,fragile]
  \frametitle{}
  \centering
  \textbf{Chapter 1}
  \centering
  Motivations for our Work
\end{frame}

\begin{frame}[fragile]
\frametitle{Problems Facing Modern Compute}
  \begin{figure}
    \centering
    \begin{minipage}[b]{0.4\textwidth}
        \centering
        \includegraphics[width=\textwidth]{images/mooreslaw.png}
        \caption{Moores Law Slowdown}
    \end{minipage}
    \hfill
    \begin{minipage}[b]{0.4\textwidth}
        \centering
        \includegraphics[width=\textwidth]{images/dennardscaling.png}
        \caption{End of Dennard Scaling}
    \end{minipage}
    \hfill
    \begin{minipage}[b]{0.5\textwidth}
        \centering
        \includegraphics[width=\textwidth]{images/vonneumann.png}
        \caption{Von-Neumann Bottleneck}
    \end{minipage}
\end{figure}
\end{frame}

\begin{frame}[fragile]
\frametitle{Overview of Modern Compute}
\framesubtitle{}
\begin{itemize}
\item Processor trade-offs: ASICs (high upfront costs), CPUs (performance limitations), GPUs (power-hungry, less flexible)
\item Question: Should modern compute be limited to CPUs/GPUs and select ASICs?
\item Four scenarios where current compute falls short:
  \begin{enumerate}
    \item Many peripherals with real-time compute needs
    \item Unusual number representations (e.g., neural network quantization)
    \item Novel architectures for existing problems (e.g., KANs)
    \item Power-efficient yet flexible solutions
  \end{enumerate}
  \end{itemize}
\end{frame}

\begin{frame}[fragile]
\frametitle{How to deal with this?}
  \framesubtitle{}
  \begin{itemize}
    \item The free lunch afforded by hardware improvements over years
      is coming to an end.
    \item New, creative, first-principles architectures and
      compute infrastructure need to be designed along with the
      software abstractions to use them. 
  \end{itemize}
\end{frame}

\begin{frame}[c,fragile]
  \frametitle{}

  \centering
  \textbf{Chapter 2}
  \centering
  Introduction To Reconfigurable And Heterogeneous Computing
\end{frame}

\begin{frame}[fragile]
  \frametitle{Setting the stage}
  \framesubtitle{}
Key concepts:
  \begin{enumerate}
    \item Reconfiguration: Reprogramming to implement new circuits (e.g., FPGAs)
    \item Heterogeneity: Integrating diverse processors (CPUs/GPUs/DSPs/ASICs) to
  \textcolor{green}{complement} rather than \textcolor{red}{replace} existing compute
  \end{enumerate}

Challenges:
  \begin{enumerate}
    \item FPGA integration with traditional software is cumbersome
    \item Custom FPGA hardware development has a steep learning curve
  \end{enumerate}
\end{frame}

\begin{frame}[fragile]
\frametitle{Reconfigurability: An Introduction to FPGAs}
\framesubtitle{}
  \begin{enumerate}
    \item Digital circuits are made up of gates and connections between
      them.
    \item FPGAs are a grid of programmable cells and switches that allows
      a user to create new gates and connections after the IC has been
      fabricated.
    \item Circuits for FPGAs are described using Hardware Descriptions Languages
      (HDLs) such as Verilog, VHDL.
    \item High level description of a circuit is compiled into real hardware
      (i.e. a representation that only uses FPGA primitives) by a "compiler".
  \end{enumerate}

\end{frame}

\begin{frame}[fragile]
\frametitle{Problem 1: Programming model for FPGAs}
\framesubtitle{}
  A true programming model for FPGAs would heavily exploit reconfigurability.

Proposal: Non-Von Neumann, flow-based reconfigurable architecture
  \begin{itemize}
    \item No instructions or ISA
    \item Data flows through the chip, transformed by configured hardware
  \end{itemize}
\end{frame}

\begin{frame}[fragile]
  \frametitle{Comparison with a Von-Neumann Computer}
  \framesubtitle{}
  \begin{figure}
    \centering
    \includegraphics[width=1.0\textwidth]{images/flow.png}
    \caption{a) A Von-Neumann Computer b) A Flowing Reconfigurable Computer}
    \label{}
  \end{figure}
\end{frame}

\begin{frame}[fragile]
  \frametitle{An exemplary DSL for reconfigurable compute architectures}
\framesubtitle{}
  \begin{lstlisting}
    Base *input = new TensorSend("ts");
    Base *proc_one = new JPEGEncoderTillDct(input, "jpeg_encoder");
    Base *proc_two = new MlEngineCore(proc_one, "ml_core");
    Base *output = new TensorReceive(proc_two, "tr");
    Model m1 = new Model(input, output);
    m1->compute(input);
  \end{lstlisting}
\end{frame}

\begin{frame}[fragile]
  \frametitle{An exemplary DSL for reconfigurable compute architectures (2)}
\framesubtitle{Continue}
   \begin{figure}
        \centering
        \includegraphics[width=1\linewidth]{images/flowcnnjpeg.png}
        \caption{Data flow for JPEG (Partial) + CNN inference}
    \end{figure}
\end{frame}


\begin{frame}[c,fragile]
  \frametitle{}

  \centering
  \textbf{Chapter 3} 
  \centering
  Need for modern EDA Compilers
\end{frame}

\newcommand\myheading[1]{%
  \par\bigskip
  {\Large\bfseries#1}\par\smallskip}

 
\begin{frame}[fragile]
  \frametitle{Problem 2: Writing Hardware Is Hard}
  \framesubtitle{}
Current state:
  \begin{itemize}
\item Proprietary tools and architectures
\item NP-Hard and NP-Complete problems in the flow
\item Gap between CAD research and production-ready tools
  \end{itemize}
\end{frame}

\begin{frame}[fragile]
  \frametitle{FPGA CAD Toolflow}
  \framesubtitle{}
   \begin{figure}
        \centering
        \includegraphics[width=1\linewidth]{images/cad_flow.png}
        \caption{FPGA CAD Tool Flow}
        \label{exa_cadflow}
    \end{figure}
\end{frame}

\begin{frame}[fragile]
  \frametitle{FPGA CAD Toolflow}
  \framesubtitle{}
   \begin{itemize}
     \item Logical Synthesis: HDL to primitive gates (AND/OR).
     \item Technology Mapping: Mapping generic AND/OR gates to k-LUTs.
     \item Placement: Placing k-LUTs in a constrained grid representing the FPGA
       based on a cost model
     \item Routing: Connection placed LUTs so that the latency is minimum
     \item Bitstream Generation: Generating a final bitstream that the FPGA
       can understand.
    \end{itemize}
\end{frame}

\begin{frame}[fragile]
  \frametitle{Optimization Opportunities for EDA Compilers}
  \framesubtitle{}
Optimization opportunities:
  \begin{itemize}
\item Direct netlist operations in DSL compiler
\item Caching optimzations for faster incremental compilation
\item GPU-accelerated routing
  \end{itemize}
\end{frame}

\begin{frame}[c,fragile]

  \centering
  \textbf{Chapter 4}
  Work Done Towards Implementation
\end{frame}

\begin{frame}[fragile]
  \frametitle{Realizing this goal}
  Vaaman: A reconfigurable heterogeneous single board computer (SBC)
  \begin{enumerate}
\item Custom-designed hardware
\item Integrates FPGA for reconfigurability
\item Heterogeneous as it integrates CPU with FPGA
  \end{enumerate}
  \framesubtitle{}
\end{frame}

\begin{frame}[fragile]
  \frametitle{The Hardware (Vaaman)}
  \framesubtitle{}
  \begin{figure}
    \centering
    \includegraphics[width=0.95\textwidth]{vaaman.jpg}
    \caption{Vaaman: A heterogenous SBC \\ \url{https://shorturl.at/5y9QA}}
    \label{neuron}
  \end{figure}
  
\end{frame}

\begin{frame}[fragile]
  \frametitle{Features}
  \begin{itemize}
    \item \textbf{Gati: CNN accelerator using FPGA as co-processor}
      \begin{enumerate}
        \item Systolic-Array based MAC engine
        \item Macro-ISA for control flow
        \item Compiler and runtime for model partitioning and execution
      \end{enumerate}

    \item \textbf{Periplex: On-the-fly peripheral generation}
      \begin{enumerate}
        \item Software-defined peripherals on FPGAs
        \item JSON-driven configuration
        \item Linux kernel driver integration
      \end{enumerate}
      \end{itemize}
\end{frame}

\begin{frame}[fragile]
  \frametitle{Conclusion}
  \framesubtitle{}
  \begin{enumerate}
    \item Reconfigurable architectures address von Neumann bottleneck and modern compute challenges
    \item Heterogeneous approach enables integration with existing infrastructure
    \item Key challenges: developing reconfigurability-focused programming models and efficient hardware compilers
    \item Solutions include novel DSLs, optimized EDA tools, and GPU acceleration
  \end{enumerate}
\end{frame}

\begin{frame}[fragile]
  \frametitle{}
  \framesubtitle{}

  Extended slides and more stuff: 

  \begin{verbatim}
  github.com/vicharak-in/noisa
  \end{verbatim}

  We are hiring! Email us at: 

  \begin{verbatim}
  join-the-team@vicharak.in
  \end{verbatim}

\end{frame}

}

\end{document}
